\chapter{Architecture Decision Records: Decision Log}
\label{app:adr-log}

This appendix summarises the most consequential Architectural Decision Records
(ADRs) for the PROTEA platform. Each record was written autonomously as
implementation evidence accumulated and confirmed by the author on 2026-05-06.
The full text of every record lives in \texttt{docs/source/adr/} inside the
PROTEA repository; this log captures the rationale and the accepted resolution
for the decisions that most directly shaped the platform's design and
operational posture.

ADRs are identified by two complementary numbering schemes. The legacy
three-digit codes (ADR-001 through ADR-008) cover the earliest architectural
choices and are referenced in Chapter~4. The \textbf{D-series} codes (D1
onward) cover the post-F2 design decisions discussed in Chapters~4 and~5 and
throughout this log. Both series are cross-referenced where they overlap.

% ────────────────────────────────────────────────────────────
\section{D4: API Versioning Strategy}
\label{app:adr-d4}

\paragraph{Context.}
PROTEA exposes a REST API consumed by the front-end, by LAFA containers, and
by downstream pipelines. As the API surface stabilised toward the F4 release,
breaking changes required a versioning strategy that would not strand existing
consumers.

\paragraph{Decision accepted 2026-05-06.}
A universal \texttt{/v1/} path prefix is applied to all routers. Future
breaking changes branch under \texttt{/v2/} without removing \texttt{/v1/}
endpoints. \texttt{Accept}-header negotiation is deferred unless a concrete
need surfaces. Deprecated unprefixed mounts are kept for one release cycle to
avoid breaking live clients (the \texttt{protea.ngrok.app} deployment and
existing front-end fetchers) during the transition.

\paragraph{Consequences.}
Each API version ships its own OpenAPI document. Schemathesis runs against the
version under test. Front-end and external clients update their base URLs once,
at the \texttt{/v1/} transition, and are thereafter insulated from internal
restructuring under the same major version.

% ────────────────────────────────────────────────────────────
\section{D6: Authentication Strategy}
\label{app:adr-d6}

\paragraph{Context.}
Sensitive endpoints (job creation, dataset import, re-ranker model upload,
evaluation triggers) were unauthenticated in the pre-F5 platform. Public cloud
exposure, LAFA submission tooling, and external adopters required an
authentication layer before the F5 phase could be declared open.

\paragraph{Decision accepted 2026-05-06.}
Two complementary mechanisms are deployed:

\begin{itemize}
  \item \textbf{API key} for service-to-service calls: an \texttt{ApiKey} ORM
    table issues opaque tokens; callers present them as
    \texttt{Authorization: Bearer \textlangle{}token\textrangle{}}.
  \item \textbf{OIDC via reverse proxy} for human users: \textbf{Authentik}
    acts as the identity provider behind \texttt{oauth2-proxy}, chosen for
    its lighter operational footprint and simpler Docker-Compose setup
    relative to Keycloak.
\end{itemize}

Rate limiting is enforced via \texttt{slowapi} on a per-endpoint basis.
A CORS allowlist is maintained in the FastAPI middleware configuration.

\paragraph{Implementation status.}
Partial as of the thesis writing phase. Tasks T5.6a (API key ORM and
middleware) and T5.6b (oauth2-proxy integration) carry the outstanding work.

% ────────────────────────────────────────────────────────────
\section{D7: Observability Stack}
\label{app:adr-d7}

\paragraph{Context.}
PROTEA's pre-F-OPS observability relied on per-process log files and ad-hoc
log statements. Multi-target deployment (cloud, HPC, airgap) and external
adopters require distributed tracing, metrics with service-level objectives,
and structured log aggregation.

\paragraph{Decision accepted 2026-05-06.}
A single canonical observability stack:

\begin{itemize}
  \item \textbf{Tracing}: OpenTelemetry (OTLP exporter) instrumenting FastAPI,
    SQLAlchemy, and \texttt{pika}. The \texttt{traceparent} header is
    propagated HTTP to queue to worker so that a single prediction is
    visible end-to-end as one OTel trace.
  \item \textbf{Metrics}: Prometheus client; a \texttt{/metrics} scrape
    endpoint is exposed by the FastAPI process.
  \item \textbf{Dashboards}: Grafana, with dashboard JSON committed under
    \texttt{deploy/grafana/} so dashboards are version-controlled.
  \item \textbf{Logs}: structured JSON via \texttt{python-json-logger},
    shipped to Loki via the \texttt{loki-docker-driver} from container
    stdout (no separate Promtail sidecar in the cloud target). Loki was
    preferred over the ELK stack because it indexes labels rather than
    full text, has a smaller memory footprint, and integrates with the
    same Grafana instance already surfacing Prometheus dashboards.
\end{itemize}

Three service-level objectives are documented in \texttt{docs/SLOs.md};
alert rules and per-alert runbooks are committed to the repository.

\paragraph{Implementation status.}
Tasks T5.1a (OTel instrumentation), T5.1b (Prometheus endpoint), T5.2
(Grafana dashboards), T5.3 (Loki driver), and T5.4 (alert rules) carry the
phased rollout; T5.1a and T5.1b are complete.

% ────────────────────────────────────────────────────────────
\section{D10: \texttt{schema\_sha} v2 Parallel Migration}
\label{app:adr-d10}

\paragraph{Context.}
\texttt{schema\_sha} is the load-bearing fingerprint that prevents inference
from running with a re-ranker booster trained against a different feature
schema. Two independent definitions of \texttt{compute\_schema\_sha} (one in
the lab, one in PROTEA) silently drifted until the parity bug surfaced during
the v9 study on 2026-05-01 and caused a non-reproducible run.

\paragraph{Decision accepted 2026-05-06.}
A parallel \texttt{schema\_sha\_v2} column is added to both \texttt{Dataset}
and \texttt{RerankerModel}. The column is backfilled from the canonical
\texttt{protea\_contracts.compute\_schema\_sha} implementation. The production
read path switches to \texttt{v2}; the \texttt{v1} column is retained until
the F3 gate for audit purposes and then dropped. A regression test exposes the
v1/v2 drift on historical rows rather than retroactively modifying them.

\paragraph{Implementation gate.}
Task T1.6 carries the Alembic migration and backfill. The migration is
classified \texttt{requires\_human}: it must be rehearsed in staging (or a
local-DB dry-run) and the backfill verified before any production rollout.
This constraint was explicitly confirmed by the author on 2026-05-06.

% ────────────────────────────────────────────────────────────
\section{D15: \texttt{protea-method} Distribution Channels}
\label{app:adr-d15}

\paragraph{Context.}
\texttt{protea-method} is the pure inference path (KNN, feature computation,
re-ranker application). Because it has no FastAPI or SQLAlchemy dependencies
it can ship independently of the full platform, addressing audiences that want
the annotation method without operating the database and messaging
infrastructure.

\paragraph{Decision accepted.}
Three distribution channels:

\begin{enumerate}
  \item \textbf{PyPI} public package (\texttt{pip install protea-method}).
    The default install is torch-free; the gated-attention MIL head is
    opt-in via the \texttt{[mil]} Poetry extra (see also D\,no\,extra below).
  \item \textbf{Docker Hub} minimal image (\texttt{protea-method-runtime},
    approximately 3 to 4\,GB with one canonical PLM and one default booster;
    CLI: \texttt{protea-predict input.fasta output.tsv}).
  \item \textbf{Companion engineering paper} covering the F-OPS pillars
    (containers, multi-target deployment, observability, secrets management).
\end{enumerate}

LAFA submission containers (D23) build on top of
\texttt{protea-method-runtime}, making D15 a prerequisite for the submission
tooling.

\paragraph{Distribution discipline.}
The \textbf{opt-in \texttt{[mil]} extra} decision (recorded separately in the
implementation log) keeps the default \texttt{pip install protea-method}
torch-free: the gated-attention MIL head is only installed when the caller
explicitly requests it. This avoids a several-gigabyte PyTorch dependency for
users who only need the LightGBM re-ranker inference path.

% ────────────────────────────────────────────────────────────
\section{D25: HPC Operation Mode}
\label{app:adr-d25}

\paragraph{Context.}
PROTEA must support HPC environments such as the BSC MareNostrum cluster.
HPC sites typically forbid privileged Docker, may restrict outbound network
access, and schedule workloads via SLURM. Two operation modes were
considered:

\begin{description}
  \item[Mode B] Stateless PROTEA workers running on HPC nodes connect to a
    PostgreSQL and RabbitMQ hosted in the cloud (LifeWatch or EOSC). No
    persistent state on the HPC node itself.
  \item[Mode C] Fully airgapped batch bundle: an Apptainer \texttt{.sif}
    image with a snapshot database pre-loaded, the default booster, and a
    single-node SLURM job requiring no outbound traffic.
\end{description}

\paragraph{Decision accepted 2026-05-06.}
\textbf{Mode B is the primary HPC deployment target.} It is closer to the
cloud architecture, reuses the same Postgres and RabbitMQ services, and
requires no special airgap provisioning. Mode C (the airgapped \texttt{.sif}
bundle, T-OPS.9) is \textbf{deferred per D25}: it becomes a post-defense
item (F9), contingent on concrete engagement with a restricted-access HPC
site. Mode B is sufficient for the thesis defense scope.

% ────────────────────────────────────────────────────────────
\section{D27: Image Registry}
\label{app:adr-d27}

\paragraph{Context.}
Seven OCI images must be hosted in a registry visible from cloud deployments,
from HPC tooling that pulls images before converting them to \texttt{.sif}
(Apptainer), and from external adopters consuming
\texttt{protea-method-runtime}.

\paragraph{Decision accepted 2026-05-06.}
\texttt{ghcr.io} (GitHub Container Registry) is the canonical registry.
GitHub Actions push images on SemVer tag using the repository's own
\texttt{GITHUB\_TOKEN}; no external registry credentials are required.
\texttt{protea-method-runtime} is published with public visibility; internal
images are org-scoped or private. Mirroring to Docker Hub is deferred until
external pull rates demand it. Task T-OPS.8 implements the \texttt{docker.yml}
workflow and is complete.

% ────────────────────────────────────────────────────────────
\section{D28: Secrets Management}
\label{app:adr-d28}

\paragraph{Context.}
PostgreSQL credentials, MinIO keys, OIDC client secrets, optional external API
tokens, and SSH keys cannot live in plaintext in any repository.
Multi-target deployment (cloud, HPC, airgap) requires a single mechanism that
works across all environments without per-environment custom solutions.

\paragraph{Decision accepted 2026-05-06.}
\texttt{sops} with \texttt{age} keys. Encrypted
\texttt{secrets.\{dev,prod\}.enc.yaml} files are committed to the repository;
CI decrypts with the age key stored as a \texttt{SOPS\_AGE\_KEY} GitHub
repository secret. Two reasons capture the choice: (a) age keys are
post-PGP ed25519, short, and revocation-chain-free; (b) sops is
file-format-agnostic so the same workflow handles YAML, JSON, and
\texttt{.env} files uniformly.

\paragraph{Implementation status.}
The scaffolding was committed alongside this ADR: \texttt{.sops.yaml},
\texttt{secrets/secrets.dev.example.yaml} (plaintext schema template,
tracked), and updated \texttt{.gitignore} rules that forbid the plaintext
\texttt{secrets.\{dev,prod\}.yaml} while keeping encrypted forms and
the example template tracked. The manual bootstrap step (generating an
\texttt{age} keypair and populating the recipient placeholder in
\texttt{.sops.yaml}) is classified \texttt{requires\_human} as task T-OPS.7.

% ────────────────────────────────────────────────────────────
\section{D29: Release Pipeline}
\label{app:adr-d29}

\paragraph{Context.}
Seven repositories need independent SemVer release cadences. Releasing one
repository must not break another; cross-repository integration testing is
required on every tag. \texttt{protea-contracts} is the most disruptive
package: version bumps propagate through all seven consumers.

\paragraph{Decision accepted 2026-05-06.}
Per-repository SemVer plus a cross-repository integration test on tag:

\begin{itemize}
  \item A SemVer tag (\texttt{vX.Y.Z}) in any repository dispatches a build,
    an image push (D27), and an integration test that pulls the new image
    together with the pinned versions of the other six repositories and runs
    a smoke pipeline.
  \item Failures block image promotion; the tag remains but the image is
    held as a release candidate until the integration test passes.
  \item A \texttt{protea-contracts} release triggers an automated re-pin PR
    in all consumers.
\end{itemize}

\texttt{release-please} is used for version bumps and CHANGELOG generation,
driven by Conventional Commits: \texttt{feat:} increments the minor version,
\texttt{fix:} increments the patch, and a \texttt{BREAKING CHANGE:} footer or
\texttt{feat!:} increments the major version. This differs from the originally
considered \texttt{python-semantic-release}, which push-commits the version
bump directly to \texttt{main} and would bypass branch protection.
\texttt{release-please} instead opens a release PR, preserving the hard project
rule that nothing lands on \texttt{main} without a passing PR review.

The commit-message convention (Conventional Commits) was already enforced from
the F2 phase onward, so the CHANGELOG generation requires no retroactive
reformatting.
